% ============================================================================
% Fichier  : partie5/P05_C17_threat_intelligence.tex
% Titre    : Le Renseignement sur la Menace
% Auteur   : MINKA MI NGUIDJOÏ Thierry Emmanuel
% ============================================================================
% Sources : MITRE ATT&CK ; Kill Chain (Lockheed Martin, 2011) ; Diamond Model
% (2013) ; STIX/TAXII (OASIS). Opération MAYO.
% ============================================================================

\chapter{Le Renseignement sur la Menace}
\label{chap:renseignement}
\chaptermark{Le renseignement sur la menace}

\epigraph{%
    « Connaître l'adversaire, c'est transformer une alerte en histoire. »%
}{Principe du cours}

\niveauChapitre{Intermédiaire à expert}{Chapitre~\ref{chap:csirt}~;
Chapitre~\ref{chap:regles}.}

\begin{boxobjectifs}
\begin{itemize}
    \item Définir le renseignement sur la menace et ses usages.
    \item Utiliser la chaîne d'attaque (Kill Chain) et ATT\&CK.
    \item Comprendre le modèle du Diamant.
    \item Situer les formats d'échange (STIX/TAXII) et les indicateurs.
    \item Mapper un incident aux techniques de l'adversaire.
\end{itemize}
\end{boxobjectifs}

\bigskip

% ============================================================================
\section*{Question directrice : l'attaque de LiZbiZ ressemble-t-elle à d'autres, et qu'en déduire~?}
% ============================================================================

Une attaque n'est presque jamais inédite~: elle suit des étapes connues et
emploie des techniques répertoriées. Le renseignement sur la menace permet de
reconnaître le mode opératoire, d'anticiper la suite et d'écrire de meilleures
règles.

\medskip
\textbf{Plan du chapitre.}
\begin{enumerate}
    \item Qu'est-ce que le renseignement (section~\ref{sec:c17-definir}).
    \item Chaîne d'attaque et ATT\&CK (section~\ref{sec:c17-attack}).
    \item Le modèle du Diamant (section~\ref{sec:c17-diamant}).
    \item Échanger le renseignement (section~\ref{sec:c17-echanger}).
\end{enumerate}

\separateur

% ============================================================================
\section{Qu'est-ce que le Renseignement}
\label{sec:c17-definir}
% ============================================================================

\begin{boxdef}
\termecle{Renseignement sur la menace}~: information analysée sur les
adversaires, leurs motivations, leurs techniques et leurs indicateurs, qui
éclaire la détection, la réponse et la décision. Il transforme des données
brutes en connaissance actionnable.
\end{boxdef}

% ============================================================================
\section{Chaîne d'Attaque et ATT\&CK}
\label{sec:c17-attack}
% ============================================================================

La \termecle{chaîne d'attaque} (Lockheed Martin, 2011) décrit une intrusion
en étapes~: reconnaissance, armement, livraison, exploitation, installation,
pilotage, actions sur objectif. \termecle{MITRE ATT\&CK} affine en tactiques
et techniques observées. Les deux se complètent~: la chaîne raconte
l'histoire, ATT\&CK en nomme les gestes.

\begin{tableaucours}{L{4.2cm}L{4.4cm}L{3.6cm}}{%
    \tableauHeader{Étape (LOGONE)} &
    \tableauHeader{Tactique ATT\&CK} &
    \tableauHeader{Trace}
}
Accès initial (compte de service) & Accès aux identifiants & Connexion anormale \\
Installation (WinUpd32) & Persistance & Service installé \\
Pilotage (connexions régulières) & Commande et contrôle & Flux réguliers \\
Exfiltration (40 Go) & Exfiltration & Volume sortant \\
Chiffrement & Impact & Fichiers chiffrés \\
\end{tableaucours}
\vspace{-0.8em}
\begin{center}{\small\itshape Tableau~17.1. L'opération LOGONE mappée à ATT\&CK~\cite{mitreattack,killchain2011}}\end{center}

% ============================================================================
\section{Le Modèle du Diamant}
\label{sec:c17-diamant}
% ============================================================================

\begin{boxdef}
\termecle{Modèle du Diamant} (2013)~: cadre d'analyse d'une intrusion par
quatre sommets reliés~: l'\emph{adversaire}, ses \emph{capacités}, son
\emph{infrastructure}, et la \emph{victime}. Relier deux sommets connus permet
d'inférer les autres.
\end{boxdef}

Exemple~: connaissant l'infrastructure de pilotage (l'adresse externe) et la
victime (LiZbiZ), on peut rapprocher l'incident d'autres victimes partageant
la même infrastructure, et ainsi caractériser l'adversaire.

% ============================================================================
\section{Échanger le Renseignement}
\label{sec:c17-echanger}
% ============================================================================

Les \termecle{indicateurs de compromission} (adresses, condensats, noms de
domaine) s'échangent entre organisations dans des formats normalisés
(\termecle{STIX} pour décrire, \termecle{TAXII} pour transporter). L'ANTIC et
les CERT sectoriels jouent ce rôle de partage au niveau national.

\begin{boxalert}
\textbf{Prudence.} Un indicateur partagé peut vieillir ou être faux.

\smallskip
\textbf{Règle~:} on qualifie la source et la fraîcheur d'un indicateur avant
d'en faire une règle de blocage~: un indicateur périmé génère du bruit, un
indicateur erroné bloque du trafic légitime.
\end{boxalert}

\begin{boxscenario}{Opération MAYO. Le mode opératoire reconnu}
\textbf{Constatation.} Mappée à ATT\&CK et au Diamant, l'attaque de LiZbiZ
correspond à un mode opératoire de rançongiciel courant en Afrique centrale~:
accès par compte de service, persistance par service, exfiltration avant
chiffrement.

\medskip
\textbf{Analyse} (confiance moyenne). La reconnaissance du mode opératoire a
permis d'anticiper la recherche de persistances cachées et d'écrire des
règles ciblées. On ne nomme pas l'adversaire avec certitude~: NEK sur
l'attribution.
\end{boxscenario}

\begin{boxmetier}{Analyste SOC / Red Team}
\textbf{Du renseignement aux règles.} L'analyste transforme chaque
caractérisation d'adversaire en règles de détection et en points de chasse~:
le renseignement ne vaut que s'il améliore la défense.
\end{boxmetier}

\separateur

\begin{boxresume}
\textbf{Ce qu'il faut retenir de ce chapitre~:}
\begin{itemize}
    \item \textbf{Le renseignement} rend la menace actionnable.
    \item \textbf{La chaîne d'attaque} raconte, \textbf{ATT\&CK} nomme les
          gestes.
    \item \textbf{Le Diamant} relie adversaire, capacités, infrastructure,
          victime.
    \item \textbf{STIX/TAXII} échangent les indicateurs~; on qualifie leur
          fraîcheur.
    \item \textbf{L'attribution} reste souvent un NEK.
\end{itemize}
\end{boxresume}

\begin{boxinfo}
\textbf{Points de contrôle de l'élève.} Je sais~: définir le renseignement~;
mapper une attaque à la chaîne et à ATT\&CK~; appliquer le Diamant~; situer
STIX/TAXII~; transformer le renseignement en règles.
\end{boxinfo}

\bigskip

\begin{boxexo}[title={\ding{45}\quad Exercice~17.1. Chaîne \niveaubase}]
Citez trois étapes de la chaîne d'attaque.
\end{boxexo}

\begin{boxexo}[title={\ding{45}\quad Exercice~17.2. Diamant \niveaubase}]
Quels sont les quatre sommets du modèle du Diamant~?
\end{boxexo}

\begin{boxexo}[title={\ding{45}\quad Exercice~17.3. Indicateur \niveaubase}]
Donnez trois exemples d'indicateurs de compromission.
\end{boxexo}

\begin{boxexo}[title={\ding{45}\quad Exercice~17.4. Mapper \niveaumoyen}]
Mappez l'installation de \texttt{WinUpd32} à une tactique ATT\&CK.
\end{boxexo}

\begin{boxexo}[title={\ding{45}\quad Exercice~17.5. Inférer \niveaumoyen}]
Connaissant l'infrastructure et la victime, qu'infère-t-on avec le Diamant~?
\end{boxexo}

\begin{boxexo}[title={\ding{45}\quad Exercice~17.6. Indicateur périmé \niveaumoyen}]
Pourquoi qualifier la fraîcheur d'un indicateur avant d'en faire une règle~?
\end{boxexo}

\begin{boxexo}[title={\ding{45}\quad Exercice~17.7. Attribution \niveaumoyen}]
Pourquoi l'attribution reste-t-elle souvent un NEK~?
\end{boxexo}

\begin{boxexo}[title={\ding{45}\quad Exercice~17.8. STIX/TAXII \niveaumoyen}]
Quelle est la différence entre STIX et TAXII~?
\end{boxexo}

\begin{boxexo}[title={\ding{45}\quad Exercice~17.9. Mapping complet \niveauexpert}]
Dressez le mapping complet de l'opération LOGONE à ATT\&CK.
\end{boxexo}

\begin{boxexo}[title={\ding{45}\quad Exercice~17.10. Du renseignement aux règles \niveauexpert}]
À partir du mode opératoire reconnu, proposez trois règles et deux points de
chasse.
\end{boxexo}

\bigskip

\begin{boxinfo}
\textbf{Pour aller plus loin.}
\begin{itemize}
    \item MITRE ATT\&CK~\cite{mitreattack}~; chaîne d'attaque~\cite{killchain2011}~;
          modèle du Diamant~\cite{diamond2013}~; STIX/TAXII~\cite{stixtaxii}.
\end{itemize}
\end{boxinfo}

\begin{boxinfo}
\textbf{Transition.} Reste à agir concrètement sur l'incident~: le
Chapitre~\ref{chap:confinement} traite du confinement, de l'éradication et du
rétablissement.
\end{boxinfo}
